% Provenance (never printed): SNT "Les reseaux sociaux" (graphs, distance, diameter, radius, centre, small worlds)
% and "Localisation, cartographie et mobilite" (positioning, road maps as graphs, NMEA frames, routes).
% See sources/book2/BRIEFS.md.
\chapter{Graphs: Networks of Roads and People}\label{ch:g10:graphs-first-look}

A rumour starts in Brenmoor: the library will open on Sundays. Each person who hears it tells their friends the next
day. If it starts with the right person, how many days until everyone who can hear it has heard it? And who is the
right person?

The question is not about anyone's personality: it is about the shape of the network of friendships. The same shape
appears in the roads between places, the links between web pages and the cables between computers. This chapter
names that shape, a \emph{\omterm{def:g10:graphs-first-look:graph}{graph}}, stores it in Python, measures \omterm{def:g10:graphs-first-look:distance}{distances} in it, and finds its \omterm{def:g10:graphs-first-look:eccentricity}{centre}.

\begin{recall}
A dictionary maps keys to values, and a set holds values without order or repetition (Book 1, grade 8, and the
chapter on tuples and sets). The Earth and space sciences book explains how a satellite receiver works out its own
position on the Earth (grade 9); a map then turns that position into a route.
\end{recall}

\section{Vertices and edges}\label{sec:g10:graphs-first-look:graphs}

\begin{definition}[Graph, vertex, edge]\label{def:g10:graphs-first-look:graph}
A \emph{graph}\index{graph} is a set of \emph{vertices}\index{vertex} together with a set of
\emph{edges}\index{edge}, each edge joining two different vertices. Two vertices joined by an edge are
\emph{adjacent}.
\end{definition}

In the friendship \omterm{def:g10:graphs-first-look:graph}{graph} of Brenmoor the \omterm{def:g10:graphs-first-look:graph}{vertices} are the people $1$ to $40$, and an \omterm{def:g10:graphs-first-look:graph}{edge} joins two friends. In the
road \omterm{def:g10:graphs-first-look:graph}{graph} the \omterm{def:g10:graphs-first-look:graph}{vertices} are the $30$ places, and an \omterm{def:g10:graphs-first-look:graph}{edge} is a road.

\begin{definition}[Degree, neighbour]\label{def:g10:graphs-first-look:degree}
The \emph{neighbours}\index{neighbour} of a \omterm{def:g10:graphs-first-look:graph}{vertex} are the \omterm{def:g10:graphs-first-look:graph}{vertices} adjacent to it. Its
\emph{degree}\index{degree} is the number of its neighbours.
\end{definition}

\begin{definition}[Directed graph, weighted graph]\label{def:g10:graphs-first-look:directed}
In a \emph{directed graph}\index{directed graph} each \omterm{def:g10:graphs-first-look:graph}{edge} goes from one \omterm{def:g10:graphs-first-look:graph}{vertex} to another, like a one-way street: an
\omterm{def:g10:graphs-first-look:graph}{edge} from $a$ to $b$ does not let you go from $b$ to $a$. A \omterm{def:g10:graphs-first-look:graph}{vertex} then has an \emph{\omterm{def:g10:graphs-first-look:degree}{out-degree}}, the number of \omterm{def:g10:graphs-first-look:graph}{edges}
leaving it, and an \emph{\omterm{def:g10:graphs-first-look:degree}{in-degree}}, the number of \omterm{def:g10:graphs-first-look:graph}{edges} arriving. In a \emph{weighted graph}\index{weighted graph}
each \omterm{def:g10:graphs-first-look:graph}{edge} carries a number, its \emph{weight}: a \omterm{def:g10:graphs-first-look:path}{length} in metres, a time, a price.
\end{definition}

\begin{omfigure}
\begin{tikzpicture}[font=\footnotesize]
  \begin{scope}
    \node[gnode] (a) at (0,1.2) {A}; \node[gnode] (b) at (1.6,1.2) {B};
    \node[gnode] (c) at (0,0) {C}; \node[gnode] (d) at (1.6,0) {D};
    \draw[gedge] (a) -- (b) (a) -- (c) (b) -- (c) (b) -- (d);
    \node at (0.8,-0.8) {undirected};
    \node[font=\scriptsize, text=omDef, anchor=west] at (1.95,1.2) {degree 3};
  \end{scope}
  \begin{scope}[xshift=45mm]
    \node[gnode] (a) at (0,1.2) {A}; \node[gnode] (b) at (1.6,1.2) {B};
    \node[gnode] (c) at (0,0) {C}; \node[gnode] (d) at (1.6,0) {D};
    \draw[ptr, black] (a) -- (b); \draw[ptr, black] (c) -- (a); \draw[ptr, black] (b) -- (c);
    \draw[ptr, black] (b) -- (d);
    \node at (0.8,-0.8) {directed};
    \node[font=\scriptsize, text=omDef, align=left, anchor=west] at (2.0,1.2) {B: in 1, out 2};
  \end{scope}
  \begin{scope}[xshift=90mm]
    \node[gnode] (a) at (0,1.2) {A}; \node[gnode] (b) at (1.6,1.2) {B};
    \node[gnode] (c) at (0,0) {C}; \node[gnode] (d) at (1.6,0) {D};
    \draw[gedge] (a) -- node[wedge, above] {4} (b) (a) -- node[wedge, left] {2} (c)
      (b) -- node[wedge] {1} (c) (b) -- node[wedge, right] {5} (d);
    \node at (0.8,-0.8) {weighted};
  \end{scope}
\end{tikzpicture}

{\small The same four \omterm{def:g10:graphs-first-look:graph}{vertices} and four \omterm{def:g10:graphs-first-look:graph}{edges}, as an undirected \omterm{def:g10:graphs-first-look:graph}{graph}, a \omterm{def:g10:graphs-first-look:directed}{directed graph} and a \omterm{def:g10:graphs-first-look:directed}{weighted graph}.}
\end{omfigure}

\section{A graph in Python}\label{sec:g10:graphs-first-look:python}

A \omterm{def:g10:graphs-first-look:graph}{graph} is stored as a dictionary from each \omterm{def:g10:graphs-first-look:graph}{vertex} to the set of its \omterm{def:g10:graphs-first-look:degree}{neighbours}. An undirected \omterm{def:g10:graphs-first-look:graph}{edge} is written at
both of its ends; a \omterm{def:g10:graphs-first-look:graph}{vertex} without any \omterm{def:g10:graphs-first-look:graph}{edge} still gets an empty set, so that it is not forgotten.

\omcode{code/grade-10/09-graphs-first-look/listings/load_friends.py}{1}{19}{The friendship \omterm{def:g10:graphs-first-look:graph}{graph} as a dictionary of
sets.}\label{lst:g10:graphs-first-look:load-friends}

\omoutput{code/grade-10/09-graphs-first-look/outputs/load_friends.txt}{Person $40$ is a \omterm{def:g10:graphs-first-look:graph}{vertex} with no
\omterm{def:g10:graphs-first-look:graph}{edge}.}\label{out:g10:graphs-first-look:load-friends}

\omcode{code/grade-10/09-graphs-first-look/listings/degree_table.py}{1}{11}{How many people have each number of
friends.}\label{lst:g10:graphs-first-look:degree-table}

\omoutput{code/grade-10/09-graphs-first-look/outputs/degree_table.txt}{The \omterm{def:g10:graphs-first-look:degree}{degrees} add up to $100$, twice the $50$
friendships.}\label{out:g10:graphs-first-look:degree-table}

\begin{proposition}[Sum of the degrees]\label{prop:g10:graphs-first-look:handshake}
In an undirected \omterm{def:g10:graphs-first-look:graph}{graph}, the sum of the \omterm{def:g10:graphs-first-look:degree}{degrees} of the \omterm{def:g10:graphs-first-look:graph}{vertices} is twice the number of \omterm{def:g10:graphs-first-look:graph}{edges}.
\end{proposition}

The proof is \cref{exo:g10:graphs-first-look:12}. A road \omterm{def:g10:graphs-first-look:graph}{graph} needs weights, and one-way roads make it directed: the
listing stores a dictionary of \omterm{def:g10:graphs-first-look:degree}{neighbours} with their \omterm{def:g10:graphs-first-look:distance}{distances}, and writes a two-way road as two directed \omterm{def:g10:graphs-first-look:graph}{edges}.

\omcode{code/grade-10/09-graphs-first-look/listings/load_roads.py}{1}{21}{The road \omterm{def:g10:graphs-first-look:graph}{graph}, directed and
weighted.}\label{lst:g10:graphs-first-look:load-roads}

\omoutput{code/grade-10/09-graphs-first-look/outputs/load_roads.txt}{$53$ roads give $105$ directed \omterm{def:g10:graphs-first-look:graph}{edges}: one road,
from the Town Square to the Library, is one-way.}\label{out:g10:graphs-first-look:load-roads}

A program that builds a \omterm{def:g10:graphs-first-look:graph}{graph} by hand can forget one end of an \omterm{def:g10:graphs-first-look:graph}{edge}. The check is short: every \omterm{def:g10:graphs-first-look:graph}{edge} must have its
reverse, except the one-way roads.

\omcode{code/grade-10/09-graphs-first-look/listings/check_symmetric.py}{1}{11}{Finding the \omterm{def:g10:graphs-first-look:graph}{edges} without a
reverse.}\label{lst:g10:graphs-first-look:check-symmetric}

\omoutput{code/grade-10/09-graphs-first-look/outputs/check_symmetric.txt}{The friendships are symmetric; among the
roads, only road $1$ is not.}\label{out:g10:graphs-first-look:check-symmetric}

\section{Paths and distance}\label{sec:g10:graphs-first-look:paths}

\begin{definition}[Path, length of a path]\label{def:g10:graphs-first-look:path}
A \emph{path}\index{path (graph)} is a sequence of \omterm{def:g10:graphs-first-look:graph}{vertices} in which each one is adjacent to the next (in a \omterm{def:g10:graphs-first-look:directed}{directed
graph}, joined by an \omterm{def:g10:graphs-first-look:graph}{edge} going from it to the next). Its \emph{length}\index{length of a path} is its number of
\omterm{def:g10:graphs-first-look:graph}{edges}, or, in a \omterm{def:g10:graphs-first-look:directed}{weighted graph}, the sum of the weights of its \omterm{def:g10:graphs-first-look:graph}{edges}.
\end{definition}

\begin{definition}[Distance]\label{def:g10:graphs-first-look:distance}
The \emph{distance}\index{distance (graph)} from a \omterm{def:g10:graphs-first-look:graph}{vertex} $a$ to a \omterm{def:g10:graphs-first-look:graph}{vertex} $b$ is the \omterm{def:g10:graphs-first-look:path}{length} of a shortest \omterm{def:g10:graphs-first-look:path}{path} from
$a$ to $b$. If there is no such \omterm{def:g10:graphs-first-look:path}{path}, there is no distance.
\end{definition}

\begin{definition}[Connected graph, component]\label{def:g10:graphs-first-look:connected}
An undirected \omterm{def:g10:graphs-first-look:graph}{graph} is \emph{connected}\index{connected graph} if there is a \omterm{def:g10:graphs-first-look:path}{path} between any two of its \omterm{def:g10:graphs-first-look:graph}{vertices}. A
\emph{connected component}\index{connected component} is a largest set of \omterm{def:g10:graphs-first-look:graph}{vertices} joined to one another by \omterm{def:g10:graphs-first-look:path}{paths}.
\end{definition}

Computing \omterm{def:g10:graphs-first-look:distance}{distances} in a large \omterm{def:g10:graphs-first-look:graph}{graph} is an algorithm of its own, written in grade 12. Here two functions are given,
in the file \code{given_graphs.py} of the grade's folder, and used without reading them: \code{distances(g, a)}
returns the dictionary of the \omterm{def:g10:graphs-first-look:distance}{distances}, in \omterm{def:g10:graphs-first-look:graph}{edges}, from \code{a} to every \omterm{def:g10:graphs-first-look:graph}{vertex} it can reach, and
\code{components(g)} returns the components, largest first.

\omcode{code/grade-10/09-graphs-first-look/listings/components.py}{1}{8}{The components of the friendship
\omterm{def:g10:graphs-first-look:graph}{graph}.}\label{lst:g10:graphs-first-look:components}

\omoutput{code/grade-10/09-graphs-first-look/outputs/components.txt}{Three components: $29$ people, $10$ people,
and person $40$ alone.}\label{out:g10:graphs-first-look:components}

\begin{omfigure}
\begin{tikzpicture}[scale=0.82, transform shape]
  % written by figures/friends_layout.py; do not edit
\draw[gedge] (0.0,2.3) -- (2.242,-0.512);
\draw[gedge] (0.0,2.3) -- (0.998,-2.072);
\draw[gedge] (0.0,2.3) -- (-1.798,-1.434);
\draw[gedge] (0.0,2.3) -- (-2.242,0.512);
\draw[gedge, omProp] (0.0,2.3) -- (3.6,-2.6);
\draw[gedge] (0.998,2.072) -- (1.798,1.434);
\draw[gedge] (0.998,2.072) -- (-1.798,-1.434);
\draw[gedge] (0.998,2.072) -- (-0.998,2.072);
\draw[gedge] (1.798,1.434) -- (1.798,-1.434);
\draw[gedge] (2.242,0.512) -- (-0.998,-2.072);
\draw[gedge] (2.242,-0.512) -- (-2.242,0.512);
\draw[gedge, omProp] (2.242,-0.512) -- (1.6,-4.6);
\draw[gedge] (1.798,-1.434) -- (-1.798,-1.434);
\draw[gedge] (1.798,-1.434) -- (-1.798,1.434);
\draw[gedge] (0.998,-2.072) -- (0.0,-2.3);
\draw[gedge] (0.0,-2.3) -- (-2.242,-0.512);
\draw[gedge] (0.0,-2.3) -- (-0.998,2.072);
\draw[gedge] (-0.998,-2.072) -- (-2.242,0.512);
\draw[gedge] (-0.998,-2.072) -- (-0.998,2.072);
\draw[gedge] (-1.798,-1.434) -- (-2.242,-0.512);
\draw[gedge] (-2.242,-0.512) -- (-2.242,0.512);
\draw[gedge] (-2.242,-0.512) -- (-1.798,1.434);
\draw[gedge, omProp] (-2.242,0.512) -- (8.593,-1.572);
\draw[gedge] (7.2,2.1) -- (6.224,1.859);
\draw[gedge] (8.176,1.859) -- (9.164,-0.745);
\draw[gedge] (8.176,1.859) -- (7.703,-2.039);
\draw[gedge] (8.176,1.859) -- (6.224,1.859);
\draw[gedge] (8.928,1.193) -- (9.164,-0.745);
\draw[gedge] (8.928,1.193) -- (5.807,-1.572);
\draw[gedge] (9.285,0.253) -- (8.593,-1.572);
\draw[gedge] (9.164,-0.745) -- (5.807,-1.572);
\draw[gedge] (6.697,-2.039) -- (6.224,1.859);
\draw[gedge] (5.807,-1.572) -- (5.236,-0.745);
\draw[gedge] (5.807,-1.572) -- (5.472,1.193);
\draw[gedge] (5.807,-1.572) -- (6.224,1.859);
\draw[gedge, omProp] (5.115,0.253) -- (1.6,-4.6);
\draw[gedge] (3.6,-2.6) -- (3.6,-6.6);
\draw[gedge] (3.6,-2.6) -- (1.868,-5.6);
\draw[gedge] (3.6,-2.6) -- (1.6,-4.6);
\draw[gedge] (3.6,-2.6) -- (2.6,-2.868);
\draw[gedge] (4.6,-2.868) -- (5.6,-4.6);
\draw[gedge] (4.6,-2.868) -- (5.332,-5.6);
\draw[gedge] (4.6,-2.868) -- (3.6,-6.6);
\draw[gedge] (5.332,-3.6) -- (1.868,-3.6);
\draw[gedge] (5.332,-3.6) -- (2.6,-2.868);
\draw[gedge] (5.332,-5.6) -- (1.868,-5.6);
\draw[gedge] (5.332,-5.6) -- (1.6,-4.6);
\draw[gedge] (5.332,-5.6) -- (1.868,-3.6);
\draw[gedge] (4.6,-6.332) -- (2.6,-6.332);
\draw[gedge] (4.6,-6.332) -- (2.6,-2.868);
\node[gnode, minimum size=5mm, font=\scriptsize, fill=omDef!25] at (0.0,2.3) {1};
\node[gnode, minimum size=5mm, font=\scriptsize, fill=white] at (0.998,2.072) {2};
\node[gnode, minimum size=5mm, font=\scriptsize, fill=white] at (1.798,1.434) {3};
\node[gnode, minimum size=5mm, font=\scriptsize, fill=white] at (2.242,0.512) {4};
\node[gnode, minimum size=5mm, font=\scriptsize, fill=white] at (2.242,-0.512) {5};
\node[gnode, minimum size=5mm, font=\scriptsize, fill=white] at (1.798,-1.434) {6};
\node[gnode, minimum size=5mm, font=\scriptsize, fill=white] at (0.998,-2.072) {7};
\node[gnode, minimum size=5mm, font=\scriptsize, fill=white] at (0.0,-2.3) {8};
\node[gnode, minimum size=5mm, font=\scriptsize, fill=white] at (-0.998,-2.072) {9};
\node[gnode, minimum size=5mm, font=\scriptsize, fill=white] at (-1.798,-1.434) {10};
\node[gnode, minimum size=5mm, font=\scriptsize, fill=white] at (-2.242,-0.512) {11};
\node[gnode, minimum size=5mm, font=\scriptsize, fill=white] at (-2.242,0.512) {12};
\node[gnode, minimum size=5mm, font=\scriptsize, fill=white] at (-1.798,1.434) {13};
\node[gnode, minimum size=5mm, font=\scriptsize, fill=white] at (-0.998,2.072) {14};
\node[gnode, minimum size=5mm, font=\scriptsize, fill=white] at (7.2,2.1) {15};
\node[gnode, minimum size=5mm, font=\scriptsize, fill=white] at (8.176,1.859) {16};
\node[gnode, minimum size=5mm, font=\scriptsize, fill=white] at (8.928,1.193) {17};
\node[gnode, minimum size=5mm, font=\scriptsize, fill=white] at (9.285,0.253) {18};
\node[gnode, minimum size=5mm, font=\scriptsize, fill=white] at (9.164,-0.745) {19};
\node[gnode, minimum size=5mm, font=\scriptsize, fill=white] at (8.593,-1.572) {20};
\node[gnode, minimum size=5mm, font=\scriptsize, fill=white] at (7.703,-2.039) {21};
\node[gnode, minimum size=5mm, font=\scriptsize, fill=white] at (6.697,-2.039) {22};
\node[gnode, minimum size=5mm, font=\scriptsize, fill=white] at (5.807,-1.572) {23};
\node[gnode, minimum size=5mm, font=\scriptsize, fill=white] at (5.236,-0.745) {24};
\node[gnode, minimum size=5mm, font=\scriptsize, fill=white] at (5.115,0.253) {25};
\node[gnode, minimum size=5mm, font=\scriptsize, fill=white] at (5.472,1.193) {26};
\node[gnode, minimum size=5mm, font=\scriptsize, fill=white] at (6.224,1.859) {27};
\node[gnode, minimum size=5mm, font=\scriptsize, fill=omDef!25] at (3.6,-2.6) {28};
\node[gnode, minimum size=5mm, font=\scriptsize, fill=white] at (4.6,-2.868) {29};
\node[gnode, minimum size=5mm, font=\scriptsize, fill=white] at (5.332,-3.6) {30};
\node[gnode, minimum size=5mm, font=\scriptsize, fill=white] at (5.6,-4.6) {31};
\node[gnode, minimum size=5mm, font=\scriptsize, fill=white] at (5.332,-5.6) {32};
\node[gnode, minimum size=5mm, font=\scriptsize, fill=white] at (4.6,-6.332) {33};
\node[gnode, minimum size=5mm, font=\scriptsize, fill=white] at (3.6,-6.6) {34};
\node[gnode, minimum size=5mm, font=\scriptsize, fill=white] at (2.6,-6.332) {35};
\node[gnode, minimum size=5mm, font=\scriptsize, fill=white] at (1.868,-5.6) {36};
\node[gnode, minimum size=5mm, font=\scriptsize, fill=white] at (1.6,-4.6) {37};
\node[gnode, minimum size=5mm, font=\scriptsize, fill=white] at (1.868,-3.6) {38};
\node[gnode, minimum size=5mm, font=\scriptsize, fill=white] at (2.6,-2.868) {39};
\node[gnode, minimum size=5mm, font=\scriptsize, fill=white] at (8.6,-4.6) {40};

\end{tikzpicture}

{\small The friendship \omterm{def:g10:graphs-first-look:graph}{graph} of people $1$ to $40$, drawn in three groups by number. The orange \omterm{def:g10:graphs-first-look:graph}{edges} join two
groups. Persons $1$ and $28$, shaded, are the \omterm{def:g10:graphs-first-look:eccentricity}{centre} of the largest component; the group on the right is a
component of its own, except $18$, $20$ and $25$, whose friends are elsewhere.}
\end{omfigure}

On a small weighted map, a shortest route can be found by hand with a \emph{frontier}: keep the \omterm{def:g10:graphs-first-look:distance}{distance} from the
start of every place reached so far, and repeatedly settle the unsettled place with the smallest one, then look at
the roads leaving it.

\begin{omfigure}
\begin{tikzpicture}[font=\footnotesize, x=11mm, y=11mm]
  \node[gnode, fill=omDef!20] (3) at (0,0) {3};
  \node[gnode] (24) at (2.4,-1.2) {24}; \node[gnode] (26) at (2.4,1.4) {26};
  \node[gnode] (1) at (5,0) {1}; \node[gnode, fill=omProp!25] (23) at (7.6,0) {23};
  \node[gnode] (16) at (4.4,-2.6) {16}; \node[gnode] (21) at (4.4,2.8) {21};
  \draw[gedge] (3) -- node[wedge] {1040} (24);
  \draw[gedge] (3) -- node[wedge] {1790} (26);
  \draw[gedge] (24) -- node[wedge] {1590} (1);
  \draw[gedge] (26) -- node[wedge] {1560} (1);
  \draw[gedge] (1) -- node[wedge] {1210} (23);
  \draw[gedge] (24) -- node[wedge] {1390} (16);
  \draw[gedge] (26) -- node[wedge] {1360} (21);
  \draw[ptr, black] (1) -- node[wedge, above=1pt] {1360, one-way} (3);
  % the route is a band BEHIND the roads: drawn last and on top, it ran through the lengths it passes
  \begin{scope}[on background layer]
    \draw[omProp!45, line width=4.5pt, line cap=round, line join=round]
      (3.center) -- (24.center) -- (1.center) -- (23.center);
  \end{scope}
  \node[anchor=west, align=left] at (8.3,0.9) {settled in order:\\3 (0), 24 (1040),\\26 (1790), 16 (2430),\\
    1 (2630), 21 (3150),\\23 (3840)};
\end{tikzpicture}

{\small Part of the town's road map, \omterm{def:g10:graphs-first-look:path}{lengths} in metres, from the Library (3) to the North Bridge (23). The route
3--24--1--23 is $3840$~m. The road from the Town Square (1) to the Library is one-way, so the way back,
23--1--3, is only $2570$~m.}
\end{omfigure}

\section{Radius, diameter and centre}\label{sec:g10:graphs-first-look:centre}

\begin{definition}[Eccentricity, radius, diameter, centre]\label{def:g10:graphs-first-look:eccentricity}
In a \omterm{def:g10:graphs-first-look:connected}{connected graph}, the \emph{eccentricity}\index{eccentricity} of a \omterm{def:g10:graphs-first-look:graph}{vertex} is its largest \omterm{def:g10:graphs-first-look:distance}{distance} to another
\omterm{def:g10:graphs-first-look:graph}{vertex}. The \emph{radius}\index{radius (graph)} of the \omterm{def:g10:graphs-first-look:graph}{graph} is the smallest eccentricity, its
\emph{diameter}\index{diameter (graph)} the largest one, and its \emph{centre}\index{centre (graph)} the set of the
\omterm{def:g10:graphs-first-look:graph}{vertices} whose eccentricity is the radius.
\end{definition}

\begin{omfigure}
\begin{tikzpicture}[font=\footnotesize, x=14mm, y=12mm]
  \foreach \v/\x/\y/\e in {A/0/0/4, B/1/0/3, C/2/0/2, D/3/0/3, E/4/0/4, F/2/-1/3, G/2/-2/4}
    \node[gnode, label={[text=omDef]above right:\e}] (\v) at (\x,\y) {\v};
  \node[gnode, fill=omProp!30] at (C) {C};
  \draw[gedge] (A) -- (B) -- (C) -- (D) -- (E) (C) -- (F) -- (G);
\end{tikzpicture}

{\small Each \omterm{def:g10:graphs-first-look:graph}{vertex} with its \omterm{def:g10:graphs-first-look:eccentricity}{eccentricity}. The \omterm{def:g10:graphs-first-look:eccentricity}{radius} is $2$, the \omterm{def:g10:graphs-first-look:eccentricity}{diameter} $4$, and the \omterm{def:g10:graphs-first-look:eccentricity}{centre} is C alone: from C,
every \omterm{def:g10:graphs-first-look:graph}{vertex} is at most two \omterm{def:g10:graphs-first-look:graph}{edges} away.}
\end{omfigure}

The rumour of the opening question is a question about \omterm{def:g10:graphs-first-look:eccentricity}{eccentricity}: started by a person of \omterm{def:g10:graphs-first-look:eccentricity}{eccentricity} $e$, it
reaches everyone in that person's component after $e$ days. On the largest component of the friendship \omterm{def:g10:graphs-first-look:graph}{graph}:

\omcode{code/grade-10/09-graphs-first-look/listings/distances_given.py}{1}{15}{\omterm{def:g10:graphs-first-look:distance}{Distances} from person $1$, then the
\omterm{def:g10:graphs-first-look:eccentricity}{radius}, \omterm{def:g10:graphs-first-look:eccentricity}{diameter} and \omterm{def:g10:graphs-first-look:eccentricity}{centre}.}\label{lst:g10:graphs-first-look:distances-given}

\omoutput{code/grade-10/09-graphs-first-look/outputs/distances_given.txt}{Persons $1$ and $28$ reach the whole
component in $4$ days; from the worst place it takes $7$.}\label{out:g10:graphs-first-look:distances-given}

\begin{remark}
$29$ people, each with at most $5$ friends, and yet everyone is within $7$ steps of everyone in the component. In
large networks of people the \omterm{def:g10:graphs-first-look:distance}{distances} are surprisingly small, a property called the \emph{small world}: even with a few
friends each, the number of people within $k$ steps multiplies at each step, so it grows very fast with $k$.
\end{remark}

\begin{method}[Finding the centre of a small graph by hand]\label{met:g10:graphs-first-look:centre}
\begin{enumerate}
  \item For each \omterm{def:g10:graphs-first-look:graph}{vertex}, write the \omterm{def:g10:graphs-first-look:distance}{distance} to every other \omterm{def:g10:graphs-first-look:graph}{vertex}, by layers: its \omterm{def:g10:graphs-first-look:degree}{neighbours} at $1$, their new
    \omterm{def:g10:graphs-first-look:degree}{neighbours} at $2$, and so on.
  \item Its \omterm{def:g10:graphs-first-look:eccentricity}{eccentricity} is the number of the last layer.
  \item The \omterm{def:g10:graphs-first-look:eccentricity}{radius} is the smallest \omterm{def:g10:graphs-first-look:eccentricity}{eccentricity}, the \omterm{def:g10:graphs-first-look:eccentricity}{diameter} the largest; the \omterm{def:g10:graphs-first-look:eccentricity}{centre} is the \omterm{def:g10:graphs-first-look:graph}{vertices} of
    \omterm{def:g10:graphs-first-look:eccentricity}{eccentricity} equal to the \omterm{def:g10:graphs-first-look:eccentricity}{radius}.
\end{enumerate}
\end{method}

A map application turns a satellite position into a \omterm{def:g10:graphs-first-look:graph}{vertex} of the road \omterm{def:g10:graphs-first-look:graph}{graph}, and a route into a shortest \omterm{def:g10:graphs-first-look:path}{path}. The
receiver gives its position as text, in a standard sentence format: \cref{exo:g10:graphs-first-look:11} decodes one.

\section{Exercises}\label{sec:g10:graphs-first-look:exercises}

\begin{exercise}[$\star$]\label{exo:g10:graphs-first-look:1}
A \omterm{def:g10:graphs-first-look:graph}{graph} has $5$ \omterm{def:g10:graphs-first-look:graph}{vertices} with \omterm{def:g10:graphs-first-look:degree}{degrees} $3$, $2$, $2$, $2$ and $1$. How many \omterm{def:g10:graphs-first-look:graph}{edges} does it have? Can a \omterm{def:g10:graphs-first-look:graph}{graph} have
\omterm{def:g10:graphs-first-look:degree}{degrees} $3$, $3$, $1$?
\end{exercise}

\begin{exercise}[$\star$]\label{exo:g10:graphs-first-look:2}
Write the dictionary of sets of the undirected \omterm{def:g10:graphs-first-look:graph}{graph} of the first drawing of the chapter (\omterm{def:g10:graphs-first-look:graph}{vertices} A, B, C, D).
\end{exercise}

\begin{exercise}[$\star$]\label{exo:g10:graphs-first-look:3}
In the \omterm{def:g10:graphs-first-look:graph}{graph} of eccentricities, give the \omterm{def:g10:graphs-first-look:distance}{distance} from A to G and from E to F. Give two different shortest \omterm{def:g10:graphs-first-look:path}{paths}
from B to D, or explain why there is only one.
\end{exercise}

\begin{exercise}[$\star$]\label{exo:g10:graphs-first-look:4}
A \omterm{def:g10:graphs-first-look:graph}{graph} has the \omterm{def:g10:graphs-first-look:graph}{edges} $1$--$2$, $2$--$3$, $4$--$5$ and the \omterm{def:g10:graphs-first-look:graph}{vertices} $1$ to $6$. List its components. How many \omterm{def:g10:graphs-first-look:graph}{edges}
would make it \omterm{def:g10:graphs-first-look:connected}{connected}, at least?
\end{exercise}

\begin{exercise}[$\star$]\label{exo:g10:graphs-first-look:5}
Write \code{degree(g, v)}, and use it to print the people of the friendship \omterm{def:g10:graphs-first-look:graph}{graph} whose \omterm{def:g10:graphs-first-look:degree}{degree} is $1$.
\end{exercise}

\begin{exercise}[$\star\star$]\label{exo:g10:graphs-first-look:6}
A function \code{add_friendship(g, a, b)} only writes \code{g[a].add(b)}. After building the \omterm{def:g10:graphs-first-look:graph}{graph} with it, which
listing of the chapter detects the problem, and what does it print? Correct the function.
\end{exercise}

\begin{exercise}[$\star\star$]\label{exo:g10:graphs-first-look:7}
Write \code{is_path(g, seq)}, which says whether a list of \omterm{def:g10:graphs-first-look:graph}{vertices} is a \omterm{def:g10:graphs-first-look:path}{path} of \code{g}. Check that
\code{[3, 2, 10, 1, 28]} is a \omterm{def:g10:graphs-first-look:path}{path} of the friendship \omterm{def:g10:graphs-first-look:graph}{graph} and \code{[3, 1]} is not.
\end{exercise}

\begin{exercise}[$\star\star$]\label{exo:g10:graphs-first-look:8}
Draw the \omterm{def:g10:graphs-first-look:graph}{graph} with \omterm{def:g10:graphs-first-look:graph}{edges} $1$--$2$, $2$--$3$, $3$--$4$, $4$--$1$, $1$--$5$. Give the \omterm{def:g10:graphs-first-look:eccentricity}{eccentricity} of each \omterm{def:g10:graphs-first-look:graph}{vertex},
the \omterm{def:g10:graphs-first-look:eccentricity}{radius}, the \omterm{def:g10:graphs-first-look:eccentricity}{diameter} and the \omterm{def:g10:graphs-first-look:eccentricity}{centre}.
\end{exercise}

\begin{exercise}[$\star\star$]\label{exo:g10:graphs-first-look:9}
In the directed drawing of the chapter, give the \omterm{def:g10:graphs-first-look:degree}{in-degree} and \omterm{def:g10:graphs-first-look:degree}{out-degree} of every \omterm{def:g10:graphs-first-look:graph}{vertex}. Check that the \omterm{def:g10:graphs-first-look:degree}{in-degrees}
and the \omterm{def:g10:graphs-first-look:degree}{out-degrees} have the same sum, and say why that is always so.
\end{exercise}

\begin{exercise}[$\star\star$]\label{exo:g10:graphs-first-look:10}
Write \code{edges_from_dict(g)}, which returns the set of the \omterm{def:g10:graphs-first-look:graph}{edges} of an undirected \omterm{def:g10:graphs-first-look:graph}{graph}, each written once as a
pair (smaller, larger). Check the sum of the \omterm{def:g10:graphs-first-look:degree}{degrees} on the friendship \omterm{def:g10:graphs-first-look:graph}{graph}.
\end{exercise}

\begin{exercise}[$\star\star$]\label{exo:g10:graphs-first-look:11}
% ledger: g10:nmea-gga
A satellite receiver sends lines of text such as

\smallskip\noindent\hspace*{1em}\texttt{\$GPGGA,093015.00,4807.0381,N,00131.5240,W,1,08,0.9,62.4,M,47.2,M,,*41}

\smallskip\noindent The fields after \texttt{GPGGA} are the time (hours, minutes, seconds), the latitude as degrees
and minutes (\texttt{ddmm.mmmm}), \texttt{N} or \texttt{S}, the longitude as \texttt{dddmm.mmmm}, \texttt{E} or
\texttt{W}, then the quality of the fix, the number of satellites, a precision figure and the altitude in metres. Give the time, and the latitude
and longitude in decimal degrees, west and south counted negative. The two digits after \texttt{*} are a check: the
exclusive or of the codes of all characters between \texttt{\$} and \texttt{*}, in hexadecimal. What is it for?
\end{exercise}

\begin{exercise}[$\star\star\star$]\label{exo:g10:graphs-first-look:12}
Prove that in an undirected \omterm{def:g10:graphs-first-look:graph}{graph} the sum of the \omterm{def:g10:graphs-first-look:degree}{degrees} is twice the number of \omterm{def:g10:graphs-first-look:graph}{edges}. Deduce that the number of
\omterm{def:g10:graphs-first-look:graph}{vertices} of odd \omterm{def:g10:graphs-first-look:degree}{degree} is even.
\end{exercise}

\begin{exercise}[$\star\star\star$]\label{exo:g10:graphs-first-look:13}
Prove that in a \omterm{def:g10:graphs-first-look:connected}{connected graph}, $\text{radius} \le \text{diameter} \le 2 \times \text{radius}$. For the second
inequality, take a \omterm{def:g10:graphs-first-look:graph}{vertex} $c$ of the \omterm{def:g10:graphs-first-look:eccentricity}{centre} and two \omterm{def:g10:graphs-first-look:graph}{vertices} $a$ and $b$, and go from $a$ to $b$ through $c$.
\end{exercise}

\begin{exercise}[$\star\star\star$]\label{exo:g10:graphs-first-look:14}
Write \code{common_friends(g, a, b)} with a set intersection, then \code{triangles(g)}, the number of groups of three
people who are all friends with one another. Why does adding the common friends of the two ends of every \omterm{def:g10:graphs-first-look:graph}{edge} count
each triangle three times? How many triangles has the friendship \omterm{def:g10:graphs-first-look:graph}{graph}?
\end{exercise}

\begin{exercise}[$\star\star\star$]\label{exo:g10:graphs-first-look:15}
A tiny web has five pages, P1 to P5, and links P1$\to$P2, P1$\to$P3, P2$\to$P3, P4$\to$P3, P5$\to$P4, P3$\to$P1.
Rank the pages by \omterm{def:g10:graphs-first-look:degree}{in-degree}. Then give each page one vote, which it shares equally among the pages it links to, and
add the votes each page receives. Which page is first in each ranking? Explain why a search engine would rather
count a link from an important page as worth more.
\end{exercise}

\section{Problem: Handshakes across Brenmoor}\label{sec:g10:graphs-first-look:problem}

\begin{problem}[{Weekend problem --- how far apart the people of the largest group are}]\label{pb:g10:graphs-first-look:1}
The program \code{exercises/handshakes.py} uses the given functions to compute the \omterm{def:g10:graphs-first-look:eccentricity}{radius}, the \omterm{def:g10:graphs-first-look:eccentricity}{diameter} and the
\omterm{def:g10:graphs-first-look:eccentricity}{centre} of the largest component of the friendship \omterm{def:g10:graphs-first-look:graph}{graph}, then again with one \omterm{def:g10:graphs-first-look:graph}{edge} removed.

\needspace{4\baselineskip}\medskip\noindent\textbf{Part I --- The graph.}
\begin{enumerate}
  \item How many \omterm{def:g10:graphs-first-look:graph}{vertices} and \omterm{def:g10:graphs-first-look:graph}{edges} has the friendship \omterm{def:g10:graphs-first-look:graph}{graph}? Check the sum of the \omterm{def:g10:graphs-first-look:degree}{degrees}.
  \item Which people have the most friends? Which have exactly one?
  \item How many \omterm{def:g10:graphs-first-look:graph}{vertices} have an odd \omterm{def:g10:graphs-first-look:degree}{degree}? Is that consistent with \cref{exo:g10:graphs-first-look:12}?
  \item Give the \omterm{def:g10:graphs-first-look:degree}{neighbours} of person $28$.
\end{enumerate}

\needspace{4\baselineskip}\medskip\noindent\textbf{Part II --- Components.}
\begin{enumerate}[resume]
  \item List the sizes of the components.
  \item Can the rumour, started by person $1$, reach person $23$? Why?
  \item Which four \omterm{def:g10:graphs-first-look:graph}{edges} join two of the three groups of numbers? Which of them join two components?
  \item How many \omterm{def:g10:graphs-first-look:graph}{edges}, at least, would make the whole \omterm{def:g10:graphs-first-look:graph}{graph} \omterm{def:g10:graphs-first-look:connected}{connected}?
\end{enumerate}

\needspace{4\baselineskip}\medskip\noindent\textbf{Part III --- Eccentricities.}
\begin{enumerate}[resume]
  \item Using the output of the \omterm{def:g10:graphs-first-look:distance}{distance} listing, give the \omterm{def:g10:graphs-first-look:eccentricity}{eccentricity} of person $1$.
  \item Find two people at \omterm{def:g10:graphs-first-look:distance}{distance} $7$ from each other, and a shortest \omterm{def:g10:graphs-first-look:path}{path} between them.
  \item Which people are at \omterm{def:g10:graphs-first-look:distance}{distance} $4$ from person $1$?
  \item How many days does the rumour take from person $3$?
\end{enumerate}

\needspace{4\baselineskip}\medskip\noindent\textbf{Part IV --- The centre, and a missing bridge.}
\begin{enumerate}[resume]
  \item Give the \omterm{def:g10:graphs-first-look:eccentricity}{radius}, the \omterm{def:g10:graphs-first-look:eccentricity}{diameter} and the \omterm{def:g10:graphs-first-look:eccentricity}{centre} of the largest component. Check
    $\text{radius} \le \text{diameter} \le 2 \times \text{radius}$.
  \item Run the program. What happens to the components, the \omterm{def:g10:graphs-first-look:eccentricity}{radius} and the \omterm{def:g10:graphs-first-look:eccentricity}{diameter} when the friendship $1$--$28$
    ends?
  \item Why does the \omterm{def:g10:graphs-first-look:eccentricity}{centre} move to $5$ and $37$?
  \item Why does removing $5$--$37$ change nothing, while removing $1$--$28$ does?
  \item What does removing $12$--$20$ do?
  \item If the library could ask one person to spread its news, whom should it ask, and how many days would the news
    take to reach the $29$ people?
  \item State the result: the \omterm{def:g10:graphs-first-look:eccentricity}{diameter} of the largest component, its \omterm{def:g10:graphs-first-look:eccentricity}{centre}, and what the loss of the friendship
    between the two central people does to the \omterm{def:g10:graphs-first-look:eccentricity}{diameter}.
\end{enumerate}
\end{problem}
